\documentclass[10pt,letterpaper]{article}
\usepackage[utf8]{inputenc}\usepackage[T1]{fontenc}\usepackage{lmodern}
\usepackage[spanish,es-noshorthands]{babel}\usepackage{csquotes}
\usepackage[top=1.55cm,bottom=0.9cm,left=1.0cm,right=1.0cm,headheight=0.5cm,headsep=0.25cm,nofoot]{geometry}
\usepackage{needspace,multicol,amsmath,amssymb,graphicx,xcolor,array,enumitem,booktabs,fancyhdr,lastpage}
\usepackage{pgfplots}\pgfplotsset{compat=1.17}\usetikzlibrary{arrows.meta}
\definecolor{guinda}{HTML}{6F1D46}\definecolor{tinta}{HTML}{000000}
\newcommand{\cas}{\framebox[1.0em]{\rule{0pt}{1.05ex}}}
\fancypagestyle{cont}{\fancyhf{}\renewcommand{\headrulewidth}{0.3pt}%
 \fancyhead[L]{\footnotesize\textbf{Nombre:} \rule{7.2cm}{0.3pt}\quad\textbf{Boleta:} \rule{2.4cm}{0.3pt}}%
 \fancyhead[R]{\footnotesize\textbf{Versión A} · hoja \thepage\ de \pageref{LastPage}}}
\pagestyle{cont}
\setlength{\parindent}{0pt}\setlength{\columnsep}{0.55cm}\setlength{\columnseprule}{0.3pt}

\begin{document}\fontsize{8.7}{10.0}\selectfont\color{tinta}\thispagestyle{empty}\vspace*{-0.7cm}
\begin{minipage}{0.09\textwidth}\includegraphics[height=1.15cm]{ipn-escudo.png}\end{minipage}%
\begin{minipage}{0.70\textwidth}\centering
{\color{guinda}\textbf{\large Instituto Politécnico Nacional · ESCA Unidad Santo Tomás}}\\[-1pt]
{\small Laboratorio: Gestión de datos, comportamiento y perfiles del consumidor digital}\\[1pt]
{\normalsize\textbf{Examen del 2.º parcial (Sesiones 4 a 8 e Influencers)}} \hfill {\small \textbf{Dr. Aboud Barsekh-Onji}}
\end{minipage}%
\begin{minipage}{0.19\textwidth}\centering
{\color{guinda}\fbox{\parbox[c][1.0cm][c]{1.3cm}{\centering\textbf{\scriptsize VERSIÓN}\\[-2pt]{\Huge\textbf{A}}}}}
\end{minipage}
\vspace{2pt}\hrule\vspace{3pt}
{\small\textbf{Nombre:} \rule{7.2cm}{0.3pt} \hfill \textbf{Boleta:} \rule{2.4cm}{0.3pt} \hfill \textbf{Fecha:} \rule{2.2cm}{0.3pt} \hfill \textbf{Calif.:} \rule{1.5cm}{0.3pt}}\\[3pt]
{\footnotesize\textbf{Duración: 50 min.} Contesta \textbf{en esta misma hoja}, con lápiz. Cada reactivo tiene cuatro opciones (a--d) y \textbf{puede haber de una a tres correctas}: marca con una \textbf{X} el cuadro de \textbf{todas} las que consideres correctas y deja \textbf{vacío} el de las demás. Si te equivocas, borra por completo la marca. Cada opción correcta marcada suma $1/k$ y cada incorrecta marcada resta $1/k$ ($k$ = número de correctas del reactivo); el mínimo por reactivo es 0, así que marcar todas las opciones no conviene. Cada reactivo vale $10/24$ de punto. \textbf{Escribe tu nombre en cada hoja.} Algunos reactivos piden leer figuras o tablas; los datos \enquote{hipotéticos} son casos construidos para el examen. Lee con atención cada enunciado.}
\vspace{2pt}\hrule\vspace{4pt}
{\footnotesize\noindent\fbox{\parbox{\dimexpr\linewidth-2\fboxsep-2\fboxrule}{\textbf{Ecuaciones de apoyo} (para los reactivos que las requieran). \textit{Regla aditiva:} $U_i=\sum_j w_j\,s_{ij}$; \textit{conjuntiva:} se exige $s_{ij}\geq u$ en todos los atributos; \textit{lexicográfica:} gana la mejor alternativa en el atributo prioritario y se considera empatada la que quede a $\tau$ o menos de ella. \textit{Creadores:} $ER=100\times\text{interacciones}/\text{seguidores}$;\ \ $CPE=\text{costo}/\text{interacciones}$;\ \ $CPM_{seg}=1000\times\text{costo}/\text{seguidores}$.}}}\vspace{3pt}
\begin{multicols}{2}
\begin{minipage}[t]{\columnwidth}\noindent\textbf{1.} Un asistente de compras con IA gana una comisión por cada venta que cierra para el comercio. Al evaluarlo con una matriz de criterios durables y con la advertencia de Mathur et al. (2019), ¿qué afirmaciones son válidas?
\vspace{1pt}\begin{enumerate}[label=\cas\ \alph*),leftmargin=2.3em,itemsep=1.2pt,topsep=1pt,parsep=0pt,partopsep=0pt]
\item Modelo de negocio: la comisión alinea sus incentivos con cerrar la venta, no necesariamente con el mejor interés del consumidor.
\item Una capa con suficiente control sobre la interfaz puede pasar de facilitar la decisión a sesgarla, sin cambiar de tecnología.
\item Intermediación de datos: la pregunta durable es quién controla la capa entre el consumidor y el punto de venta, y quién es dueño del dato de comportamiento.
\item Por ser un asistente de IA y no una plataforma, el criterio de modelo de negocio deja de aplicar y solo importa su exposición regulatoria.
\end{enumerate}\end{minipage}\par\vspace{3pt}
\begin{minipage}[t]{\columnwidth}\noindent\textbf{2.} Se contrastan dos predicciones razonables: el fin de las cookies de terceros en Chrome y una caída de 25\,\% de la búsqueda tradicional para 2026 (Gartner, 2024). ¿Qué lecturas son correctas a septiembre de 2026?
\vspace{1pt}\begin{enumerate}[label=\cas\ \alph*),leftmargin=2.3em,itemsep=1.2pt,topsep=1pt,parsep=0pt,partopsep=0pt]
\item Las fuentes de industria se contradicen sobre si la búsqueda cayó; por eso lo no confirmable en fuente primaria no se reporta como dato duro.
\item Las predicciones fallaron porque Gartner es una fuente de menor jerarquía que la literatura arbitrada; el caso es una crítica a esa firma.
\item Las cookies de terceros siguen activas por defecto en Chrome, porque Google revirtió su plan.
\item Como la caída de 25\,\% se anunció en un informe de industria de prestigio, debe reportarse como dato duro.
\end{enumerate}\end{minipage}\par\vspace{3pt}
\begin{minipage}[t]{\columnwidth}\noindent\textbf{3.} La Figura 1 muestra, durante 104 semanas, el índice de interés de búsqueda de una marca (datos sintéticos). El equipo observa el alza de las semanas 101--104 (zona sombreada), coincidente con una campaña de \textit{influencer marketing}, y quiere duplicar el presupuesto. ¿Qué afirmaciones son válidas?
\begin{center}
\begin{tikzpicture}\begin{axis}[width=8.0cm,height=3.3cm,xmin=1,xmax=104,ymin=0,ymax=100,xlabel={semana},ylabel={índice de interés},
 axis lines*=left,ymajorgrids,grid style={black!15},tick label style={font=\tiny},label style={font=\tiny},xtick={1,13,26,39,52,65,78,91,104},ytick={0,50,100},xlabel style={yshift=3pt},ylabel style={yshift=-3pt}]
\fill[black!22] (axis cs:101,0) rectangle (axis cs:104,100);
\addplot[black,line width=0.7pt,mark=none] coordinates {(1,62.8) (2,69.3) (3,72.8) (4,73.0) (5,76.3) (6,78.4) (7,82.4) (8,82.8) (9,85.1) (10,80.9) (11,86.9) (12,83.2) (13,83.3) (14,80.1) (15,77.4) (16,76.2) (17,73.7) (18,68.4) (19,64.7) (20,62.1) (21,55.0) (22,49.5) (23,48.5) (24,42.1) (25,35.4) (26,33.2) (27,29.7) (28,24.6) (29,20.6) (30,20.2) (31,19.1) (32,14.8) (33,12.1) (34,12.9) (35,12.8) (36,10.8) (37,12.7) (38,14.5) (39,13.8) (40,14.7) (41,19.3) (42,22.0) (43,26.9) (44,26.4) (45,31.5) (46,35.5) (47,38.4) (48,46.4) (49,51.8) (50,54.3) (51,62.7) (52,66.2) (53,70.1) (54,73.8) (55,78.5) (56,77.8) (57,84.2) (58,86.7) (59,84.4) (60,89.7) (61,89.6) (62,91.8) (63,89.5) (64,89.5) (65,88.9) (66,85.0) (67,85.4) (68,81.4) (69,79.4) (70,74.0) (71,73.1) (72,68.2) (73,62.5) (74,59.6) (75,56.3) (76,52.8) (77,42.3) (78,42.5) (79,37.7) (80,32.8) (81,27.7) (82,28.7) (83,21.4) (84,22.5) (85,22.0) (86,15.5) (87,17.2) (88,15.2) (89,18.4) (90,21.6) (91,24.0) (92,19.2) (93,23.9) (94,29.1) (95,34.0) (96,35.7) (97,37.6) (98,43.8) (99,47.8) (100,52.1) (101,55.8) (102,62.5) (103,67.0) (104,70.8)};
\end{axis}\end{tikzpicture}\\[-1pt]{\tiny\textbf{Figura 1.} Interés de búsqueda de la marca, 2 años (sintético). Sombreado: semanas 101--104.}
\end{center}
\vspace{1pt}\begin{enumerate}[label=\cas\ \alph*),leftmargin=2.3em,itemsep=1.2pt,topsep=1pt,parsep=0pt,partopsep=0pt]
\item La tendencia y la estacionalidad estimadas en las últimas semanas son tan confiables como las del centro de la serie.
\item Un promedio móvil de 52 semanas estima la tendencia porque, al abarcar un ciclo anual completo, la temporada se cancela.
\item El alza de unos 15 puntos en las semanas 101--104 es evidencia suficiente de que la campaña funciona.
\item El alza equivalente de las semanas 49--52, sin campaña, sugiere un patrón anual que podría explicar el salto.
\end{enumerate}\end{minipage}\par\vspace{3pt}
\begin{minipage}[t]{\columnwidth}\noindent\textbf{4.} Un equipo descarga 5\,000 comentarios públicos de una red social sobre su marca, extrae los 10 términos más frecuentes y presenta una nube de palabras como \enquote{análisis de sentimiento}. ¿Qué observaciones son correctas?
\vspace{1pt}\begin{enumerate}[label=\cas\ \alph*),leftmargin=2.3em,itemsep=1.2pt,topsep=1pt,parsep=0pt,partopsep=0pt]
\item La nube sirve para explorar; para interpretar hace falta clasificar los términos (positivo, fricción, neutro) con un trabajo adicional.
\item Que los comentarios sean públicos no basta para justificar cualquier uso analítico posterior; hay que verificar la LFPDPPP vigente.
\item Una palabra como \textit{precio} puede ser positiva (\enquote{precio justo}) o de fricción (\enquote{subió mucho}): la frecuencia sola no lo distingue.
\item Como el análisis es agregado y no publica nombres, la LFPDPPP deja de aplicar al corpus.
\end{enumerate}\end{minipage}\par\vspace{3pt}
\begin{minipage}[t]{\columnwidth}\noindent\textbf{5.} Respecto de lo que es Google Analytics y lo que se puede esperar de él, ¿qué afirmaciones son correctas?
\vspace{1pt}\begin{enumerate}[label=\cas\ \alph*),leftmargin=2.3em,itemsep=1.2pt,topsep=1pt,parsep=0pt,partopsep=0pt]
\item Complementa, y no sustituye, a las encuestas, la escucha social y el CRM.
\item Como es una plataforma gratuita y consolidada, sus menús y nombres son estables y un tutorial de hace varios años sigue siendo una guía confiable.
\item Es una fuente de datos de primera parte y de señal conductual: registra lo que las personas hacen en un sitio web o una app propios.
\item Permite saber por qué las personas actúan como lo hacen, por lo que puede sustituir a las encuestas y a la escucha social.
\end{enumerate}\end{minipage}\par\vspace{3pt}
\begin{minipage}[t]{\columnwidth}\noindent\textbf{6.} Una clienta entra a una tienda desde su celular y desde su laptop, sin iniciar sesión, y después lo hace desde una tablet. Respecto de cómo Google Analytics la reconoce, ¿qué afirmaciones son correctas?
\vspace{1pt}\begin{enumerate}[label=\cas\ \alph*),leftmargin=2.3em,itemsep=1.2pt,topsep=1pt,parsep=0pt,partopsep=0pt]
\item Si inicia sesión y la organización implementa User-ID, Google Analytics puede unir sus dispositivos en un solo usuario.
\item El caso ilustra el problema de la identidad única que separa lo multicanal de lo omnicanal (Verhoef et al., 2015).
\item Aparece como tres usuarios, porque el \texttt{client\_id} (cookie \texttt{\_ga}) reconoce navegadores, no personas.
\item La cookie \texttt{\_ga} caduca por defecto a las 24 horas, por lo que cada visita se cuenta como un usuario nuevo.
\end{enumerate}\end{minipage}\par\vspace{3pt}
\begin{minipage}[t]{\columnwidth}\noindent\textbf{7.} Una analista consulta una propiedad estándar de Google Analytics con retención de 14 meses. Hace una Exploración que abarca 20 meses, la compara con un informe estándar del mismo periodo y quiere exportar la tabla; sus compañeros practican con la cuenta de demostración. ¿Qué afirmaciones son correctas?
\vspace{1pt}\begin{enumerate}[label=\cas\ \alph*),leftmargin=2.3em,itemsep=1.2pt,topsep=1pt,parsep=0pt,partopsep=0pt]
\item La retención de 14 meses borra los informes estándar anteriores a ese límite.
\item La Exploración se corta en torno a los 14 meses, aunque el informe estándar agregado muestre todo el periodo.
\item El archivo exportado trae una fila por persona, por lo que sirve para identificar a cada cliente.
\item En la cuenta de demostración no se puede exportar porque es de solo consulta; es una restricción de esa cuenta, no un error del usuario.
\end{enumerate}\end{minipage}\par\vspace{3pt}
\begin{minipage}[t]{\columnwidth}\noindent\textbf{8.} La tabla muestra una Exploración de embudo de Google Analytics (usuarios que alcanzan cada paso), separada por tipo de dispositivo (datos hipotéticos). ¿Qué afirmaciones son correctas?
\begin{center}{\scriptsize\renewcommand{\arraystretch}{1.05}\begin{tabular}{@{}lrr@{}}\toprule
Paso (evento) & Móvil & Escritorio\\\midrule
\texttt{view\_item} & 12\,000 & 8\,000\\ \texttt{add\_to\_cart} & 3\,000 & 2\,800\\
\texttt{begin\_checkout} & 1\,500 & 1\,960\\ \texttt{purchase} & 900 & 1\,372\\\bottomrule\end{tabular}}\\[1pt]{\tiny\textbf{Tabla 1.} Embudo de Google Analytics por dispositivo (usuarios).}\end{center}
\vspace{1pt}\begin{enumerate}[label=\cas\ \alph*),leftmargin=2.3em,itemsep=1.2pt,topsep=1pt,parsep=0pt,partopsep=0pt]
\item En ambos dispositivos el mayor abandono ocurre entre \texttt{view\_item} y \texttt{add\_to\_cart} (75\,\% en móvil; 65\,\% en escritorio).
\item Como el mayor abandono es el del pago, conviene rediseñar primero el \textit{checkout} en ambos dispositivos.
\item El paso \texttt{begin\_checkout}$\to$\texttt{purchase} avanza más en escritorio (70\,\%) que en móvil (60\,\%).
\item La conversión de punta a punta es 7.5\,\% en móvil y 17.15\,\% en escritorio.
\end{enumerate}\end{minipage}\par\vspace{3pt}
\begin{minipage}[t]{\columnwidth}\noindent\textbf{9.} Una tienda en línea revisa sus elementos de urgencia, escasez y prueba social. Según el criterio visto en clase (información veraz y salida tan fácil como la entrada), ¿qué afirmaciones son correctas?
\vspace{1pt}\begin{enumerate}[label=\cas\ \alph*),leftmargin=2.3em,itemsep=1.2pt,topsep=1pt,parsep=0pt,partopsep=0pt]
\item Un nudge se distingue de un patrón oscuro por el tamaño de su efecto en la conversión, no por la veracidad ni por la facilidad de salida.
\item Un aviso \enquote{quedan 2 piezas} que muestra siempre el mismo número, sin relación con el inventario, es un patrón oscuro.
\item Según Luguri y Strahilevitz (2021), los patrones leves casi cuadruplicaron la probabilidad de conservar un servicio no pedido y los agresivos más que la duplicaron.
\item Un temporizador de oferta es siempre un patrón oscuro, aunque informe la fecha real en que termina la promoción.
\end{enumerate}\end{minipage}\par\vspace{3pt}
\begin{minipage}[t]{\columnwidth}\noindent\textbf{10.} La tabla da los puntajes ya normalizados de 0 a 10 (10 = lo mejor; precio y entrega ya están invertidos) de cuatro alternativas ficticias. Pesos: precio 0.5, calidad 0.3, entrega 0.2. Reglas: aditiva; conjuntiva con umbral 5 en cada atributo; lexicográfica con prioridad en calidad (tolerancia 0.5). ¿Qué afirmaciones son correctas?
\vspace{1pt}\begin{enumerate}[label=\cas\ \alph*),leftmargin=2.3em,itemsep=1.2pt,topsep=1pt,parsep=0pt,partopsep=0pt]
\item La regla aditiva elige A (7.3), aunque tenga el puntaje de calidad más bajo.
\item A gana con la aditiva y queda fuera con la conjuntiva porque la primera deja que un atributo compense a otro y la segunda exige un mínimo en cada uno.
\item La regla lexicográfica con prioridad en calidad elige D, porque su entrega es la mejor.
\item La regla conjuntiva con umbral 5 deja pasar solo a B y a D.
\end{enumerate}\end{minipage}\par\vspace{3pt}
\begin{minipage}[t]{\columnwidth}\noindent\textbf{11.} Una persona con poco tiempo debe elegir entre 40 modelos de celular en una tienda en línea. Según Simon (1955), Payne et al. (1988) y Tversky y Kahneman (1974), ¿qué comportamientos son esperables?
\vspace{1pt}\begin{enumerate}[label=\cas\ \alph*),leftmargin=2.3em,itemsep=1.2pt,topsep=1pt,parsep=0pt,partopsep=0pt]
\item Descartará primero con una regla no compensatoria (p. ej., un mínimo por atributo) y comparará después las sobrevivientes con una compensatoria.
\item Ponderará todos los atributos de las 40 alternativas con una regla aditiva, porque la interfaz digital lo facilita.
\item Las heurísticas de representatividad, disponibilidad y anclaje no se ven afectadas por cómo la interfaz ordena y presenta las opciones.
\item Buscará una opción suficientemente buena en lugar de la óptima (racionalidad limitada).
\end{enumerate}\end{minipage}\par\vspace{3pt}
\begin{minipage}[t]{\columnwidth}\noindent\textbf{12.} Se reportan las siguientes cifras: 86.1\,\% de la población de 6 años o más usa internet y 37.3\,\% de las personas usuarias compró en línea (ENDUTIH 2025); el comercio electrónico minorista fue 17.7\,\% de las ventas minoristas y hubo 77.2 millones de compradores digitales (AMVO, 2026). ¿Qué lecturas son correctas?
\vspace{1pt}\begin{enumerate}[label=\cas\ \alph*),leftmargin=2.3em,itemsep=1.2pt,topsep=1pt,parsep=0pt,partopsep=0pt]
\item El 17.7\,\% indica que casi una de cada cinco personas usuarias de internet compra en línea.
\item Que 37.3\,\% haya comprado en línea no implica que 62.7\,\% decida fuera de internet (puede haber \textit{webrooming}), y las dos fuentes miden cosas distintas, por lo que sus cifras no se suman.
\item El \textit{webrooming} es probar en tienda y comprar en línea; por eso la cifra de ENDUTIH subestima a quienes compran en tienda.
\item Como 37.3\,\% compró en línea, 62.7\,\% decide su compra fuera de internet.
\end{enumerate}\end{minipage}\par\vspace{3pt}
\begin{minipage}[t]{\columnwidth}\noindent\textbf{13.} Sobre el origen y el respaldo de los marcos del recorrido del cliente, ¿qué afirmaciones son correctas?
\vspace{1pt}\begin{enumerate}[label=\cas\ \alph*),leftmargin=2.3em,itemsep=1.2pt,topsep=1pt,parsep=0pt,partopsep=0pt]
\item El acrónimo AIDA fue acuñado por Strong en 1925, al sistematizar la secuencia que Lewis formuló en 1898.
\item El \textit{consumer decision journey} (Court et al., 2009) es investigación de una consultora y no fue arbitrado: es un marco útil, no un experimento neutral.
\item La atribución de AIDA a Lewis (1898) circula desde 1925 y, según Iwamoto (2024), carece de evidencia primaria.
\item La crítica al embudo lineal solo la sostienen consultoras y plataformas; no tiene respaldo académico independiente.
\end{enumerate}\end{minipage}\par\vspace{3pt}
\begin{minipage}[t]{\columnwidth}\noindent\textbf{14.} Con la lectura digital de las etapas de AIDAS (atención, interés, deseo, acción, satisfacción), ¿qué pares de etapa, dato y métrica son correctos?
\vspace{1pt}\begin{enumerate}[label=\cas\ \alph*),leftmargin=2.3em,itemsep=1.2pt,topsep=1pt,parsep=0pt,partopsep=0pt]
\item Acción: el \textit{checkout} deja la transacción; la métrica es la tasa de conversión, que por ser la mejor medida suele sobrevalorarse.
\item Atención: el clic y el tiempo en página son el dato; la métrica asociada es la tasa de clics.
\item Satisfacción: la calificación y la devolución son el dato; las métricas asociadas son la recompra y la recomendación.
\item Deseo: guardar o agregar al carrito deja el evento «agregar al carrito»; la métrica asociada es el abandono de carrito.
\end{enumerate}\end{minipage}\par\vspace{3pt}
\begin{minipage}[t]{\columnwidth}\noindent\textbf{15.} Un equipo usa un modelo de recorrido tipo cadena de Markov, que calcula la probabilidad de pasar de cada etapa a la siguiente, para comparar tres escenarios (resultados hipotéticos en la tabla). ¿Qué afirmaciones son correctas?
\begin{center}\resizebox{\columnwidth}{!}{\scriptsize\renewcommand{\arraystretch}{1.05}\begin{tabular}{@{}lccc@{}}\toprule
Escenario & \% que convierte & Contactos hasta salir & Redes: último clic / remoción\\\midrule
A: base & 30 & 4.0 & 4\,\% / 20\,\%\\ B: más retrocesos & 22 & 5.5 & 3\,\% / 17\,\%\\ C: más redes, mejor directo & 36 & 3.8 & 6\,\% / 31\,\%\\\bottomrule\end{tabular}}\\[1pt]{\tiny\textbf{Tabla 2.} Resultados hipotéticos de un modelo de recorrido.}\end{center}
\vspace{1pt}\begin{enumerate}[label=\cas\ \alph*),leftmargin=2.3em,itemsep=1.2pt,topsep=1pt,parsep=0pt,partopsep=0pt]
\item Como el modelo supone que el siguiente paso depende de todo el historial de la persona, los resultados de los escenarios son exactos.
\item Comparar escenarios permite estimar el efecto de cambiar un supuesto (p. ej., más retrocesos) sin gastar presupuesto, pero depende de que las probabilidades de transición reflejen el comportamiento real.
\item El escenario B convierte menos y tarda más que el A: volver atrás a evaluar baja la conversión y la retrasa; un embudo lineal que lo ignora sobreestimaría la conversión.
\item En el escenario C el crédito de redes sociales sube con ambos modelos de atribución, por lo que el último clic y el efecto de remoción miden lo mismo.
\end{enumerate}\end{minipage}\par\vspace{3pt}
\begin{minipage}[t]{\columnwidth}\noindent\textbf{16.} Una marca reporta con el modelo de último clic: redes sociales 4\,\%, búsqueda 52\,\%, directo 38\,\%, correo 6\,\%. Un modelo de efecto de remoción estima que apagar redes sociales reduciría cerca de 20\,\% la probabilidad global de convertir. La dirección propone pasar el presupuesto de redes a búsqueda de pago. ¿Qué interpretaciones son válidas?
\vspace{1pt}\begin{enumerate}[label=\cas\ \alph*),leftmargin=2.3em,itemsep=1.2pt,topsep=1pt,parsep=0pt,partopsep=0pt]
\item El efecto de remoción es una medida del modelo, no una prueba causal; para confirmarlo convendría un experimento (p. ej., apagar el canal en una región).
\item Si los cálculos del modelo de último clic son aritméticamente correctos, no puede llevar a una mala decisión.
\item Trasladar el presupuesto a búsqueda es la decisión más segura, porque la acción es la etapa mejor medida y la más cercana a la causa de la compra.
\item Redes sociales suele abrir recorridos que otros canales cierran; el último clic no ve ese papel.
\end{enumerate}\end{minipage}\par\vspace{3pt}
\begin{minipage}[t]{\columnwidth}\noindent\textbf{17.} Diego (16 años) quiere unos audífonos para videojuegos tras ver las reseñas de un \textit{streamer} al que sigue. Su madre fija el presupuesto y compara tres reseñas; su abuelo los compra con su tarjeta. Diego los usa y publica un \textit{unboxing}. Según los siete papeles de la decisión de compra, ¿qué afirmaciones son correctas?
\vspace{1pt}\begin{enumerate}[label=\cas\ \alph*),leftmargin=2.3em,itemsep=1.2pt,topsep=1pt,parsep=0pt,partopsep=0pt]
\item El \textit{streamer} es el decisor, porque su reseña fue lo que más pesó en la compra.
\item El abuelo es el comprador, y sus datos de transacción describen al usuario del producto.
\item Diego es usuario y consumidor, y al publicar el \textit{unboxing} también es prosumidor.
\item La madre es la decisora y, al comparar reseñas, también ejerce el papel de evaluadora.
\end{enumerate}\end{minipage}\par\vspace{3pt}
\begin{minipage}[t]{\columnwidth}\noindent\textbf{18.} Una marca compara tres creadores de Instagram (datos hipotéticos de la tabla). ¿Qué afirmaciones son correctas?
\begin{center}{\scriptsize\renewcommand{\arraystretch}{1.05}\begin{tabular}{@{}lrrr@{}}\toprule
Creador & Seguidores & Interacciones & Tarifa (MXN)\\\midrule
X & 30\,000 & 1\,200 & \$6\,000\\ Y & 300\,000 & 3\,000 & \$36\,000\\ Z & 90\,000 & 9\,900 & \$18\,000\\\bottomrule\end{tabular}}\\[1pt]{\tiny\textbf{Tabla 3.} Cotizaciones por publicación (hipotéticas).}\end{center}
\vspace{1pt}\begin{enumerate}[label=\cas\ \alph*),leftmargin=2.3em,itemsep=1.2pt,topsep=1pt,parsep=0pt,partopsep=0pt]
\item Z tiene el \textit{engagement rate} más alto (11\,\%) y el menor costo por interacción (\$1.82), pero un engagement tan alto justifica auditar su audiencia antes de elegirlo.
\item Si el objetivo es notoriedad, Y es más eficiente que X por su menor costo por mil seguidores; si el objetivo es interacción, X es más eficiente que Y por su menor costo por interacción.
\item Y tiene un \textit{engagement rate} mayor que X porque suma más interacciones totales (3\,000 frente a 1\,200).
\item Un \textit{engagement rate} mayor siempre indica una audiencia de mejor calidad, por lo que Z debe elegirse.
\end{enumerate}\end{minipage}\par\vspace{3pt}
\begin{minipage}[t]{\columnwidth}\noindent\textbf{19.} Sobre la influencia de los creadores en la decisión de compra, ¿qué afirmaciones son correctas?
\vspace{1pt}\begin{enumerate}[label=\cas\ \alph*),leftmargin=2.3em,itemsep=1.2pt,topsep=1pt,parsep=0pt,partopsep=0pt]
\item La Guía de Publicidad para Influencers de PROFECO (2023) es una ley específica con régimen sancionador propio.
\item Las cifras de interacción por tamaño de audiencia (${\sim}$2.7\,\% nano; $<$1\,\% macro) provienen de un informe primario ya verificado.
\item Los metaanálisis de Leung et al. (2022) y de Barari et al. (2025) convergen en que el efecto opera por credibilidad y atractivo percibidos.
\item Un influenciador pequeño puede ser un buen punto de contacto de confianza, pero es una hipótesis que cada marca debe contrastar con sus datos.
\end{enumerate}\end{minipage}\par\vspace{3pt}
\begin{minipage}[t]{\columnwidth}\noindent\textbf{20.} Sobre los datos que deja cada figura participante y su lectura, ¿qué afirmaciones son correctas?
\vspace{1pt}\begin{enumerate}[label=\cas\ \alph*),leftmargin=2.3em,itemsep=1.2pt,topsep=1pt,parsep=0pt,partopsep=0pt]
\item Si el evaluador es un asistente de IA, deja de importar con qué regla pondera, porque la configura el usuario.
\item El 97.6\,\% de uso de internet de los jóvenes de 15 a 24 años (ENDUTIH 2025) indica a quién debe dirigirse el anuncio.
\item El decisor es la figura mejor medida, porque su elección siempre queda registrada como transacción.
\item Los datos de uso describen al usuario y los de compra al comprador; del decisor, a menudo no hay ningún dato visible.
\end{enumerate}\end{minipage}\par\vspace{3pt}
\begin{minipage}[t]{\columnwidth}\noindent\textbf{21.} Sobre cómo clasificar y entender a un creador de contenido, ¿qué afirmaciones son correctas?
\vspace{1pt}\begin{enumerate}[label=\cas\ \alph*),leftmargin=2.3em,itemsep=1.2pt,topsep=1pt,parsep=0pt,partopsep=0pt]
\item Los cortes por seguidores (nano $<$10 mil; micro 10--100 mil; macro 100 mil--1 M; mega $>$1 M) son universales.
\item La relación parasocial es de ida y vuelta: el creador conoce a cada seguidor, y por eso la recomendación se siente de confianza.
\item Una médica con 8\,000 seguidores que explica temas de salud es nano por tamaño y experta (líder de opinión) por rol.
\item Un mismo creador puede ser micro por tamaño, experto por rol y afiliado por modelo de pago.
\end{enumerate}\end{minipage}\par\vspace{3pt}
\begin{minipage}[t]{\columnwidth}\noindent\textbf{22.} Sobre los datos de engagement y de pago a creadores, ¿qué afirmaciones son correctas?
\vspace{1pt}\begin{enumerate}[label=\cas\ \alph*),leftmargin=2.3em,itemsep=1.2pt,topsep=1pt,parsep=0pt,partopsep=0pt]
\item Casi todos los \textit{benchmarks} vienen de empresas del sector y de muestras distintas, por lo que se compara dentro de la misma plataforma, nicho y mercado.
\item En el modelo de comisión o afiliado se paga una tarifa fija por publicación, independiente de las ventas.
\item En TikTok, una muestra de 1\,500 creadores mostró que el engagement no cae de forma constante con el tamaño (nano 7.0\,\%, micro 3.1\,\%, macro 3.9\,\%, mega 6.5\,\%).
\item Las tarifas por publicación en México (p. ej., \$300--\$2\,000 para nano) forman un tarifario oficial.
\end{enumerate}\end{minipage}\par\vspace{3pt}
\begin{minipage}[t]{\columnwidth}\noindent\textbf{23.} Tres creadores cotizan una publicación (datos hipotéticos). ¿Qué afirmaciones son correctas?
\vspace{1pt}\begin{enumerate}[label=\cas\ \alph*),leftmargin=2.3em,itemsep=1.2pt,topsep=1pt,parsep=0pt,partopsep=0pt]
\item Si el objetivo es conversión, B es la mejor opción por su menor CPM.
\item A y C tienen el mismo costo por interacción (\$5.00), aunque A tiene mayor \textit{engagement rate}.
\item C tiene el \textit{engagement rate} más alto (3\,\%).
\item B es el más barato por cada mil seguidores (\$120), pero tiene el \textit{engagement rate} más bajo de los tres (1\,\%).
\end{enumerate}\end{minipage}\par\vspace{3pt}
\begin{minipage}[t]{\columnwidth}\noindent\textbf{24.} Sobre los riesgos de contratar creadores y la transparencia publicitaria en México, ¿qué afirmaciones son correctas?
\vspace{1pt}\begin{enumerate}[label=\cas\ \alph*),leftmargin=2.3em,itemsep=1.2pt,topsep=1pt,parsep=0pt,partopsep=0pt]
\item Según la guía de PROFECO (2023), el contenido pagado debe identificarse como publicidad; es una guía de recomendaciones y su fuerza vinculante sigue en discusión.
\item La publicidad no declarada puede ser una infracción y, además, destruye la confianza que se pagó por tener.
\item La cifra de IAB México de 3.3 millones de creadores equivale a la inversión publicitaria oficial en influencers.
\item Un engagement muy superior a 10\,\% en Instagram y comentarios genéricos repetidos son señales de audiencia inflada (bots, \textit{pods}).
\end{enumerate}\end{minipage}\par\vspace{3pt}
\end{multicols}\end{document}